\documentclass[11pt]{article}
\usepackage{amsmath,amsthm,amssymb,bm}
\usepackage[margin=2.2cm]{geometry}
\usepackage{microtype}

\newcommand{\vv}[1]{\bm{#1}}
\newcommand{\mm}[1]{\mathbf{#1}}

\theoremstyle{plain}
\newtheorem{theorem}{Theorem}
\newtheorem{lemma}{Lemma}
\theoremstyle{definition}
\newtheorem{definition}{Definition}

\begin{document}
\begin{center}{\large\bfseries Positive Definiteness of $K_d$}\end{center}
\smallskip

\noindent\textbf{Setup.}
Let $\vv{d}\in\mathbb{R}^n$ be the Lagrangian finger-space coordinates,
$\vv{d}_{\rm ref}$ the virtual equilibrium, and
$\vv{\delta d}=\vv{d}_{\rm ref}-\vv{d}$ the current deflection.
The predicted tip force along contact normal $\vv{n}$ is
\[
  f = \vv{n}^{\top}\mm{A}\,K_d\,\vv{\delta d},
  \qquad
  \mm{A} = \bigl[(\mm{P}\mm{J}_x)^*\bigr]^{\!\top}\mm{\eta}\,\mm{J}_d^{\top}.
\]
Define $\vv{v}=\mm{A}^{\top}\vv{n}\in\mathbb{R}^n$ (the force-sensitivity direction in
finger space) and $\vv{u}=\vv{\delta d}\in\mathbb{R}^n$ (the current deflection),
so that $f=\vv{v}^{\top}K_d\,\vv{u}$.
Set $\gamma=\|\vv{v}\|^2\|\vv{u}\|^2$.

\begin{definition}
$K_d$ is \emph{positive definite} if $\vv{x}^{\top}K_d\,\vv{x}>0$ for all
$\vv{x}\neq\vv{0}$, equivalently if its symmetric part
$K_s:=\tfrac{1}{2}(K_d+K_d^{\top})$ is positive definite in the usual sense.
\end{definition}

% -------------------------------------------------------------------
\section*{1.\; Force-tracking update}

The loss $\mathcal{L}=\tfrac{1}{2}(f-f_{\rm des})^2$ has gradient
$\partial\mathcal{L}/\partial K_d = e\,\vv{v}\vv{u}^{\top}$
where $e=f-f_{\rm des}$ is the current force error.
The update rule is
\[
  K_d^{(t+1)} = K_d^{(t)} - \alpha\,e_t\,\vv{v}\vv{u}^{\top}.
\]
Note that $\vv{v}$ and $\vv{u}$ are fixed (quasi-static assumption:
configuration and reference held constant during the update).

\begin{theorem}\label{thm:main}
Let $K_d^{(0)}\succ 0$, $f_{\rm des}>0$, and $\alpha\in(0,2/\gamma)$.
Then $K_d^{(t)}\succ 0$ for all $t\geq 0$,
for every such $\alpha$, if and only if the initial force overshoot satisfies
\begin{equation}\label{eq:cond}
  e_0 = f^{(0)} - f_{\rm des}
  \;<\;
  \lambda_{\min}(K_s^{(0)})\cdot\|\vv{v}\|\,\|\vv{u}\|.
\end{equation}
\end{theorem}

\begin{proof}
\textbf{Step 1: telescoping bound.}
Let $\mm{P}=\tfrac{1}{2}(\vv{v}\vv{u}^{\top}+\vv{u}\vv{v}^{\top})$
be the symmetric part of the rank-1 update direction.
Each step changes $K_s$ by $-\alpha e_t\mm{P}$.
Since the largest eigenvalue of $\mm{P}$ satisfies
$\lambda_+(\mm{P})=(\vv{v}^{\top}\vv{u}+\|\vv{v}\|\|\vv{u}\|)/2\leq\|\vv{v}\|\|\vv{u}\|$
(Cauchy--Schwarz), Weyl's inequality gives
\[
  \lambda_{\min}(K_s^{(t+1)}) \geq \lambda_{\min}(K_s^{(t)}) - \alpha\,|e_t|\,\|\vv{v}\|\|\vv{u}\|.
\]
Telescoping over all steps up to $T$:
\[
  \lambda_{\min}(K_s^{(T)})
  \;\geq\;
  \lambda_{\min}(K_s^{(0)})
  \;-\;
  \alpha\,\|\vv{v}\|\|\vv{u}\|\sum_{t=0}^{T-1}|e_t|.
\]

\textbf{Step 2: $\alpha$ cancels.}
Since $e_t = e_0(1-\alpha\gamma)^t$ decays geometrically,
\[
  \alpha\sum_{t=0}^{\infty}|e_t|
  \;=\;
  \alpha\cdot\frac{|e_0|}{\alpha\gamma}
  \;=\;
  \frac{|e_0|}{\gamma},
\]
which is independent of $\alpha$.
Therefore, for all $T$:
\[
  \lambda_{\min}(K_s^{(T)})
  \;\geq\;
  \lambda_{\min}(K_s^{(0)}) - \frac{|e_0|}{\|\vv{v}\|\|\vv{u}\|}.
\]
This lower bound is strictly positive if and only if \eqref{eq:cond} holds,
and it applies uniformly in $T$ and in $\alpha\in(0,2/\gamma)$.

\textbf{When $e_0\leq 0$:} $K_s^{(t)}\succeq K_s^{(0)}\succ 0$ at every step.
\end{proof}

\noindent
\textit{Key observation.}
Reducing $\alpha$ does not help: a smaller learning rate produces smaller
steps but proportionally more of them. The total eigenvalue shift
$\alpha\sum_t|e_t|\cdot\|\vv{v}\|\|\vv{u}\| = |e_0|/\|\vv{v}\|\|\vv{u}\|$
is the same for every $\alpha$.
If \eqref{eq:cond} fails, $K_d$ eventually becomes indefinite
regardless of how small $\alpha$ is.

\begin{lemma}[Isotropic initialization]\label{lem:iso}
If $K_d^{(0)}=k_0\mm{I}$ with $k_0>0$, then \eqref{eq:cond} holds automatically
and $K_d^{(t)}\succ 0$ for all $t\geq 0$ and all $\alpha\in(0,2/\gamma)$.
\end{lemma}
\begin{proof}
By Cauchy--Schwarz,
$f^{(0)} = k_0\,\vv{v}^{\top}\vv{u} \leq k_0\|\vv{v}\|\|\vv{u}\|$,
so $e_0 \leq k_0\|\vv{v}\|\|\vv{u}\| - f_{\rm des}$.
Hence
\[
  \frac{e_0}{\|\vv{v}\|\|\vv{u}\|}
  \;\leq\; k_0 - \frac{f_{\rm des}}{\|\vv{v}\|\|\vv{u}\|}
  \;<\; k_0 = \lambda_{\min}(K_s^{(0)}).
\]
Condition \eqref{eq:cond} is satisfied, and Theorem~\ref{thm:main} applies.
\end{proof}

% -------------------------------------------------------------------
\section*{2.\; Feedforward stiffness inversion}

The first-order tip stiffness is the congruence transformation
$K_x=\mm{M}^{\top}K_d\,\mm{M}$ where $\mm{M}=\mm{J}_d(\mm{P}\mm{J}_x)^*$
(Theorem~1 of the paper).
Given a desired tip stiffness $K_{x,{\rm des}}\succ 0$ and invertible $\mm{M}$,
inverting yields $K_d=(\mm{M}^{\top})^{-1}K_{x,{\rm des}}\,\mm{M}^{-1}$.

\begin{theorem}
$K_d=(\mm{M}^{\top})^{-1}K_{x,{\rm des}}\,\mm{M}^{-1}$ is positive definite.
\end{theorem}
\begin{proof}
For any $\vv{x}\neq\vv{0}$, the vector $\vv{y}=\mm{M}^{-1}\vv{x}$ is also nonzero
(since $\mm{M}$ is invertible), and
$\vv{x}^{\top}K_d\,\vv{x}=\vv{y}^{\top}K_{x,{\rm des}}\,\vv{y}>0$.
\end{proof}

\end{document}
